\subsubsection{Cubic gap}

\begin{theorem}[The 7-site gap is a cubic irrational, root of $x^3-5x^2+6x-1$]\label{thm:tile7-cubic-gap}\lean{Infomaterial.Flow.tile7_gap_root_of_cos7}\leanok\uses{thm:tile5-golden-gap}
The first tile gap that escapes the square roots is the 7-site gap. The value
$2\cos(\pi/7)$ is a root of $x^3-x^2-2x+1$; the characteristic reduction
$B(2-x)=-A(x)$ carries that to the gap cubic
$B=x^3-5x^2+6x-1$ for $\lambda_1(7)=2-2\cos(\pi/7)$. Unlike the $N\le6$
cases, this is a genuine degree-3 algebraic gap.
\end{theorem}

\begin{theorem}[Intermediate-value bracket isolating the 7-site gap]\label{thm:tile7-gap-bracket}\lean{Infomaterial.Flow.tile7_gap_bracket}\leanok\uses{thm:tile7-cubic-gap}
$B(0)=-1<0$ and $B(1)=1>0$, so the cubic gap has a root in $(0,1)$, locating
the spectral gap $2(1-\cos(\pi/7))$. The matching bracket for
$2\cos(\pi/7)\in(1,2)$ is \texttt{cos7\_bracket}. This supplies the honest
constructive locator for the casus-irreducibilis value; the Bishop real by
rational bisection remains part of the general frontier.
\end{theorem}
